% Volume VIII: Computational Neuroscience Mathematics
% Continuity note for Chapter 43.
% Previous dependencies: Chapters 8 and 25 supply ODE notation and linear dynamics; Chapter 26 supplies PDE intuition; Chapters 13-15 supply stochastic and Bayesian language for later channel and spike models.
% New durable concepts: membrane voltage, capacitance, conductance, reversal potential, passive membrane equation, membrane time constant, voltage-gated channel variables, Hodgkin-Huxley equations, cable equation, length constant, boundary conditions, compartmental approximation, and neuromorphic RC links.
% Forward hooks: Chapter 44 simplifies conductance-based biophysics into point-neuron spike models; Chapter 45 turns presynaptic spikes into synaptic currents and plasticity rules; Chapter 46 treats spikes statistically as point processes; Chapters 49 and 50 reuse biophysical and circuit language in population models and predictive processing.
% Notation commitments: membrane voltage is V(t) or V(x,t), measured as inside minus outside; currents use the sign convention that positive applied current depolarizes the cell in equations of the form C dV/dt = - ionic outward currents + I_ext; conductances are g, capacitances C, reversal potentials E, and time constants tau.
\section{Chapter 43: Neuronal Biophysics and Cable Theory}
\label{ch:neuronal-biophysics-cable-theory}

The mathematical study of neurons begins with a physical fact: a cell membrane separates charge. That separation gives a voltage, the membrane stores charge like a capacitor, and ion channels let charge leak or flow selectively. Spikes, synapses, dendritic filtering, and neuromorphic circuits all build on this electrical substrate.

\paragraph{Chapter problem.}
How can a neuron be modeled as an electrical system whose voltage changes in time, whose ion channels create nonlinear feedback, and whose dendrites distribute voltage across space?

\paragraph{Section map.}
Section 43.1 derives the passive membrane equation from an RC circuit and interprets the time constant. Section 43.2 turns ion channels into conductance and gating-variable models. Section 43.3 assembles those pieces into the Hodgkin-Huxley equations for action potentials. Section 43.4 extends the point membrane to cable theory, where voltage depends on space as well as time.

\subsection{43.1 Electrical models of membranes}

\paragraph{Guiding question.}
If a neuron is treated as one electrically compact compartment, what differential equation describes its membrane voltage?

\paragraph{Beginner-facing intuition.}
The membrane is a thin insulator between two conductive fluids. An insulator between conductors is a capacitor: it stores charge when voltage changes. Ion channels are pathways through the membrane and act like conductances. If a conductance is open, current flows according to a driving force: the voltage \(V\) is pulled toward a reversal potential \(E\). The passive membrane equation is therefore just charge conservation for a capacitor with a leak.

This model is deliberately humble. It does not yet explain spikes. It says what happens when the membrane is below threshold and the channel conductances are approximately constant. That baseline matters because every more elaborate neuron model contains this passive relaxation as one limiting case.

\paragraph{Formal definitions and notation.}
Let \(V(t)\in\mathbb{R}\) be the membrane voltage at time \(t\), measured as intracellular potential minus extracellular potential, usually in millivolts. Let \(C>0\) be the membrane capacitance of the compartment, with units farads. Let \(g_L>0\) be the leak conductance, with units siemens, and let \(E_L\) be the leak reversal potential. Let \(I_{\mathrm{ext}}(t)\) be an applied or synaptic current, written positive when it depolarizes the membrane in the equation below.

The passive membrane equation is
\[
C\frac{dV}{dt}=-g_L(V-E_L)+I_{\mathrm{ext}}(t).
\]
Define the input resistance and membrane time constant by
\[
R_m=\frac{1}{g_L},\qquad \tau_m=\frac{C}{g_L}=R_mC.
\]
Dividing by \(g_L\) gives the equivalent form
\[
\tau_m\frac{dV}{dt}=-(V-E_L)+R_m I_{\mathrm{ext}}(t).
\]
The state space is one-dimensional: \(V(t)\in\mathbb{R}\). The model parameters are \(C,g_L,E_L\), and the input is the function \(I_{\mathrm{ext}}:\mathbb{R}_{\geq0}\to\mathbb{R}\).

\paragraph{Core derivation: constant-current solution.}
Assume \(I_{\mathrm{ext}}(t)=I_0\) is constant. The equation becomes
\[
\frac{dV}{dt}=-\frac{1}{\tau_m}\bigl(V-V_\infty\bigr),
\qquad
V_\infty=E_L+R_m I_0.
\]
Subtract \(V_\infty\) and set \(u(t)=V(t)-V_\infty\). Then
\[
\frac{du}{dt}=-\frac{1}{\tau_m}u,
\]
so
\[
\boxed{V(t)=V_\infty+\bigl(V(0)-V_\infty\bigr)e^{-t/\tau_m}.}
\]
Thus the voltage does not jump instantly to its steady value. It relaxes exponentially. After one time constant, the distance to the steady state has been multiplied by \(e^{-1}\approx0.37\).

\paragraph{Concrete example.}
Suppose \(C=200\) pF and \(g_L=10\) nS. Then
\[
\tau_m=\frac{200\ \mathrm{pF}}{10\ \mathrm{nS}}=20\ \mathrm{ms},
\qquad
R_m=100\ \mathrm{M}\Omega.
\]
If \(E_L=-65\) mV and \(I_0=100\) pA, then \(R_m I_0=10\) mV and \(V_\infty=-55\) mV. Starting from \(V(0)=-65\) mV,
\[
V(20\ \mathrm{ms})=-55+(-10)e^{-1}\approx -58.7\ \mathrm{mV}.
\]
The same current would have a smaller effect in a cell with larger leak conductance because \(R_m=1/g_L\) would be smaller.

\paragraph{Computational neuroscience example.}
In a current-clamp experiment, a small square pulse of current is injected into a neuron while voltage is recorded. For a passive membrane, the voltage response is a rising exponential during the pulse and a decaying exponential after the pulse. Fitting these exponentials gives estimates of \(R_m\), \(C\), and \(\tau_m\). In a compartmental model, each compartment has a local version of this RC equation before coupling to neighboring compartments is added.

\paragraph{AI and machine-learning example.}
Leaky integration is also a computational primitive. A recurrent unit of the form
\[
\tau \dot h=-h+Wx+Uh
\]
is a continuous-time analogue of a leaky membrane with input. Neuromorphic circuits implement similar equations with capacitors and conductances, so the same time-constant algebra determines memory duration, input filtering, and energy use. The analogy is mathematical, not a claim that all artificial recurrent units are biophysical neurons.

\paragraph{Common misconceptions and failure modes.}
First, the resting potential is not zero voltage; it is a reference-dependent electrical state, often near \(-70\) mV in mammalian neurons. Second, capacitance does not leak current by itself; capacitance stores charge and determines how fast voltage changes for a given net current. Third, the sign of current depends on convention. Always inspect the equation before saying whether a positive current is inward, outward, depolarizing, or hyperpolarizing. Fourth, a one-compartment model assumes spatial compactness; long dendrites violate that assumption.

\paragraph{Exercises.}
\begin{enumerate}[label=\textbf{43.1.\arabic*.}, leftmargin=*]
\item \textbf{Mechanical.} For \(C=100\) pF and \(g_L=5\) nS, compute \(R_m\) and \(\tau_m\).
\item \textbf{Proof or derivation.} Starting from \(C\,dV/dt=-g_L(V-E_L)+I_0\), derive the exponential solution for \(V(t)\) under constant current.
\item \textbf{Numerical/computational.} Let \(E_L=-70\) mV, \(R_m=80\ \mathrm{M}\Omega\), \(\tau_m=16\) ms, \(I_0=150\) pA, and \(V(0)=E_L\). Compute \(V(16\ \mathrm{ms})\).
\item \textbf{Interpretation.} A neuron has a small time constant but a large input resistance. What does each fact say about its response to a brief current pulse?
\item \textbf{Misconception/debugging.} A student says, ``If no current is injected, voltage must be zero.'' Correct the statement using \(E_L\) and the passive membrane equation.
\end{enumerate}

\subsection{43.2 Ion channels and conductance models}

\paragraph{Guiding question.}
How do ion channels turn voltage into current, and how can channel opening be modeled by state variables?

\paragraph{Beginner-facing intuition.}
A channel does two things. It selects an ion species, which fixes a reversal potential, and it opens or closes, which changes conductance. The current is large when many channels are open and when the membrane voltage is far from the reversal potential. Voltage-gated channels add feedback: voltage changes gate opening, and gate opening changes voltage.

The key modeling move is to replace many microscopic channels by a fraction \(x(t)\in[0,1]\) of gates in a particular state. That fraction evolves smoothly in a deterministic population model, even though individual channels open and close stochastically.

\paragraph{Formal definitions and notation.}
For channel type \(i\), let \(\bar g_i\geq0\) be the maximal conductance and \(E_i\) its reversal potential. Let \(p_i(V,t)\in[0,1]\) be the fraction of channels open, or an open-probability factor. The current is modeled as
\[
I_i(V,t)=\bar g_i\,p_i(V,t)(V-E_i).
\]
In the membrane equation this current is usually subtracted when it is written as an outward ionic current:
\[
C\frac{dV}{dt}=I_{\mathrm{ext}}-\sum_i \bar g_i p_i(V,t)(V-E_i).
\]

A single gating variable \(x(t)\in[0,1]\) can obey the rate equation
\[
\frac{dx}{dt}=\alpha_x(V)(1-x)-\beta_x(V)x,
\]
where \(\alpha_x(V)\geq0\) is an opening rate and \(\beta_x(V)\geq0\) is a closing rate. Equivalently,
\[
\boxed{\frac{dx}{dt}=\frac{x_\infty(V)-x}{\tau_x(V)},}
\]
where
\[
x_\infty(V)=\frac{\alpha_x(V)}{\alpha_x(V)+\beta_x(V)},
\qquad
\tau_x(V)=\frac{1}{\alpha_x(V)+\beta_x(V)}.
\]
For multiple independent gates, an open probability may have the form \(p_i=m^p h^q\), with powers counting how many gates must be permissive.

\paragraph{Core derivation: voltage-clamp relaxation.}
Under voltage clamp, \(V(t)=V_c\) is held fixed. Then \(x_\infty(V_c)\) and \(\tau_x(V_c)\) are constants, and
\[
\frac{dx}{dt}=-\frac{1}{\tau_x(V_c)}\bigl(x-x_\infty(V_c)\bigr).
\]
Solving as in Section 43.1 gives
\[
\boxed{x(t)=x_\infty(V_c)+\bigl(x(0)-x_\infty(V_c)\bigr)e^{-t/\tau_x(V_c)}.}
\]
This equation explains why voltage clamp is mathematically powerful: by holding \(V\) fixed, it breaks the voltage-gating feedback loop and turns gate dynamics into a simple exponential fit.

\paragraph{Concrete example.}
Suppose a sodium activation gate has \(x(0)=0.10\), \(x_\infty(V_c)=0.80\), and \(\tau_x(V_c)=2\) ms after a voltage step. After \(2\) ms,
\[
x(2)=0.80+(0.10-0.80)e^{-1}\approx0.542.
\]
If a channel has open factor \(x^3\), then the open factor changes from \(0.001\) at \(x=0.10\) to about \(0.159\) at \(x=0.542\). A modest-looking change in a gate variable can produce a large conductance change.

As a current example, let \(\bar g=20\) nS, \(p=0.125\), \(V=-30\) mV, and \(E_{\mathrm{Na}}=50\) mV. Then \(g=\bar g p=2.5\) nS and
\[
I=g(V-E_{\mathrm{Na}})=2.5\ \mathrm{nS}(-80\ \mathrm{mV})=-200\ \mathrm{pA}.
\]
With the sign convention above, this inward sodium current depolarizes the cell.

\paragraph{Computational neuroscience example.}
Voltage-clamp protocols estimate \(x_\infty(V)\) and \(\tau_x(V)\) by stepping the voltage to different command values and fitting current transients. A conductance-based model then uses those fitted functions inside a current-clamp simulation. This is how channel kinetics become parameters in an ordinary differential equation model instead of verbal descriptions such as ``fast sodium'' or ``slow potassium.''

\paragraph{AI and machine-learning example.}
Gates also appear in machine learning, for example in LSTM and GRU units, attention masks, and multiplicative modulation. The mathematical similarity is that a bounded variable controls information flow. The difference is important: an LSTM gate is learned to improve a task objective, while an ion-channel gate is a physical state variable with voltage-dependent transition rates. Differentiable conductance-based neuron models can connect the two by fitting channel parameters or surrogate gating functions from data.

\paragraph{Common misconceptions and failure modes.}
First, conductance and current are not the same. Conductance says how open the pathway is; current also depends on the driving force \(V-E_i\). Second, opening an excitatory channel does not always create the same current: the current shrinks as \(V\) approaches the reversal potential. Third, a gating variable is a population fraction or probability-like state, not a single channel's continuous opening amount. Fourth, voltage clamp does not describe normal operation; it is an experimental trick that isolates channel kinetics.

\paragraph{Exercises.}
\begin{enumerate}[label=\textbf{43.2.\arabic*.}, leftmargin=*]
\item \textbf{Mechanical.} Given \(\alpha(V)=3\ \mathrm{ms}^{-1}\) and \(\beta(V)=1\ \mathrm{ms}^{-1}\) at a fixed voltage, compute \(x_\infty(V)\) and \(\tau_x(V)\).
\item \textbf{Proof or derivation.} Starting from \(\dot x=\alpha(1-x)-\beta x\), derive the \(x_\infty,\tau_x\) form.
\item \textbf{Numerical/computational.} For \(\bar g=5\) nS, \(p=0.4\), \(V=-60\) mV, and \(E=-80\) mV, compute \(I=\bar g p(V-E)\). Is the current inward or outward under the displayed sign convention?
\item \textbf{Interpretation.} Why does a voltage-gated inward current create positive feedback during depolarization?
\item \textbf{Misconception/debugging.} A student says, ``If a channel is open, it always depolarizes the neuron.'' Give a counterexample using a potassium reversal potential below rest.
\end{enumerate}

\subsection{43.3 Hodgkin-Huxley equations}

\paragraph{Guiding question.}
How can voltage-dependent sodium and potassium conductances generate an all-or-none action potential?

\paragraph{Beginner-facing intuition.}
The Hodgkin-Huxley model replaces an abstract threshold with explicit conductance feedback. Fast sodium activation creates a regenerative upstroke: depolarization opens sodium channels, sodium current depolarizes the membrane further, and the process accelerates. Slower sodium inactivation and potassium activation stop the spike and bring the voltage back down. A spike is therefore a trajectory of a coupled ODE system, not a magic discontinuity.

\paragraph{Formal definitions and notation.}
The classical Hodgkin-Huxley state is
\[
(V,m,h,n)\in\mathbb{R}\times[0,1]^3.
\]
Here \(V\) is voltage, \(m\) is sodium activation, \(h\) is sodium inactivation, and \(n\) is potassium activation. With external current \(I_{\mathrm{ext}}\), the voltage equation is
\[
C\frac{dV}{dt}
=I_{\mathrm{ext}}
-\bar g_{\mathrm{Na}}m^3h(V-E_{\mathrm{Na}})
-\bar g_{\mathrm{K}}n^4(V-E_{\mathrm{K}})
-g_L(V-E_L).
\]
Each gate obeys
\[
\frac{dx}{dt}=\alpha_x(V)(1-x)-\beta_x(V)x
=\frac{x_\infty(V)-x}{\tau_x(V)},
\qquad x\in\{m,h,n\}.
\]
The model is a four-dimensional nonlinear ordinary differential equation. Its parameters include \(C\), maximal conductances, reversal potentials, and voltage-dependent rate functions.

\paragraph{Core mechanism: positive and delayed negative feedback.}
Consider a depolarizing perturbation from rest. The fast variable \(m\) rapidly increases, so \(m^3h\) rises and sodium conductance increases. Since \(E_{\mathrm{Na}}\) is far above rest, the term \(-\bar g_{\mathrm{Na}}m^3h(V-E_{\mathrm{Na}})\) is depolarizing when \(V<E_{\mathrm{Na}}\). This is positive feedback.

The feedback is not permanent. The inactivation variable \(h\) decreases during depolarization, reducing sodium availability, while \(n\) increases and opens potassium conductance. Since \(E_{\mathrm{K}}\) is below rest, potassium current repolarizes and often hyperpolarizes the membrane. After the spike, \(h\) must recover and \(n\) must fall before another spike can be triggered easily. This gives refractoriness.

\paragraph{Concrete example.}
Use the classical parameter scale
\[
C=1\ \mu\mathrm{F}/\mathrm{cm}^2,\quad
\bar g_{\mathrm{Na}}=120\ \mathrm{mS}/\mathrm{cm}^2,\quad
\bar g_{\mathrm{K}}=36\ \mathrm{mS}/\mathrm{cm}^2,\quad
g_L=0.3\ \mathrm{mS}/\mathrm{cm}^2.
\]
If \(m=0.05\), \(h=0.60\), and \(n=0.32\), then
\[
\bar g_{\mathrm{Na}}m^3h
=120(0.05)^3(0.60)\approx0.009\ \mathrm{mS}/\mathrm{cm}^2,
\]
while
\[
\bar g_{\mathrm{K}}n^4
=36(0.32)^4\approx0.377\ \mathrm{mS}/\mathrm{cm}^2.
\]
At rest, sodium activation is tiny and potassium plus leak dominate. If depolarization raises \(m\) to \(0.50\) while \(h\) is still \(0.60\), the sodium conductance becomes \(120(0.5)^3(0.6)=9\ \mathrm{mS}/\mathrm{cm}^2\), a thousand-fold increase from the resting value in this toy calculation. That sharp nonlinearity explains the rapid upstroke.

\paragraph{Computational neuroscience example.}
The Hodgkin-Huxley equations are used when spike waveform, refractory state, pharmacological channel manipulation, or energy cost matters. For example, blocking sodium channels in the model means lowering \(\bar g_{\mathrm{Na}}\); this can abolish spikes. Increasing a slow potassium conductance can lengthen the afterhyperpolarization and reduce firing rate. The variables are interpretable enough that a simulation can be compared with voltage-clamp and current-clamp experiments.

\paragraph{AI and machine-learning example.}
Hodgkin-Huxley neurons are rarely used as the default units in large artificial networks because they are expensive to simulate and their detailed spike waveform is often unnecessary for task optimization. They are useful in differentiable simulators, neuromorphic hardware validation, biologically constrained recurrent models, and surrogate-gradient studies that ask how close simplified spiking units are to conductance-based dynamics.

\paragraph{Common misconceptions and failure modes.}
First, the Hodgkin-Huxley model does not impose a spike threshold by hand; threshold-like behavior emerges from nonlinear dynamics. Second, the variables \(m,h,n\) are not voltages and not currents; they are gating state variables. Third, the powers \(m^3\) and \(n^4\) are model choices fitted to conductance kinetics, not universal laws for all channels. Fourth, matching one spike waveform does not guarantee a model will predict all firing regimes or all cell types.

\paragraph{Exercises.}
\begin{enumerate}[label=\textbf{43.3.\arabic*.}, leftmargin=*]
\item \textbf{Mechanical.} In the HH sodium current \(\bar g_{\mathrm{Na}}m^3h(V-E_{\mathrm{Na}})\), identify the maximal conductance, gating factor, and driving force.
\item \textbf{Proof or derivation.} Explain algebraically why increasing \(m\) from \(0.1\) to \(0.4\) has a much larger-than-linear effect on \(m^3h\) when \(h\) is fixed.
\item \textbf{Numerical/computational.} Compute \(m^3h\) for \(m=0.2,h=0.8\) and for \(m=0.6,h=0.4\). Which state has larger sodium availability?
\item \textbf{Interpretation.} Describe how the post-spike state \(h\) low and \(n\) high explains refractoriness.
\item \textbf{Misconception/debugging.} A simulation shows no spikes even under depolarizing current. List two HH parameters or state variables you would inspect before concluding the code is broken.
\end{enumerate}

\subsection{43.4 Cable theory}

\paragraph{Guiding question.}
When a dendrite or axon is not electrically compact, how does voltage spread and decay over space?

\paragraph{Beginner-facing intuition.}
A dendrite is not one point. Current can flow along the inside of the dendrite, leak through the membrane, and charge membrane capacitance at every location. Cable theory turns a neurite into a distributed RC circuit. The result is a partial differential equation for \(V(x,t)\), where \(x\) is position along the cable.

The two most important scales are the time constant \(\tau_m\), which controls temporal relaxation, and the length constant \(\lambda\), which controls spatial attenuation. A synapse far from the soma can have less somatic influence because passive voltage decays along the cable.

\paragraph{Formal definitions and notation.}
Let \(V(x,t)\) be voltage along a cylindrical cable, and let \(v(x,t)=V(x,t)-E_L\) be deviation from leak reversal. In a passive cable, a common nondimensionalized form is
\[
\boxed{\tau_m\frac{\partial v}{\partial t}
=\lambda^2\frac{\partial^2 v}{\partial x^2}-v+r_m i_{\mathrm{ext}}(x,t).}
\]
Here \(\tau_m=r_m c_m\) is the membrane time constant per unit area, \(\lambda>0\) is the length constant, \(r_m\) is specific membrane resistance, \(c_m\) is specific membrane capacitance, and \(i_{\mathrm{ext}}\) is external current density in compatible units. For a cylinder of radius \(a\) and axial resistivity \(r_L\),
\[
\lambda=\sqrt{\frac{a r_m}{2r_L}}.
\]

Boundary conditions complete the model. A sealed end at \(x=L\) has no axial current leaving the end, often written
\[
\frac{\partial v}{\partial x}(L,t)=0.
\]
A voltage-clamped end specifies \(v(0,t)=v_0(t)\). A current-injected boundary specifies the spatial derivative through axial Ohm's law. Without boundary and initial conditions, the PDE is not a complete prediction rule.

\paragraph{Core derivation: steady exponential attenuation.}
Set \(i_{\mathrm{ext}}=0\) away from an injection site and look for a steady-state solution, so \(\partial v/\partial t=0\). The passive cable equation becomes
\[
\lambda^2\frac{d^2v}{dx^2}-v=0.
\]
The characteristic equation is
\[
\lambda^2 r^2-1=0,
\]
so \(r=\pm1/\lambda\). Thus
\[
v(x)=A e^{-x/\lambda}+B e^{x/\lambda}.
\]
For a semi-infinite cable extending to \(x>0\) with voltage bounded as \(x\to\infty\), we set \(B=0\), giving
\[
\boxed{v(x)=v(0)e^{-x/\lambda}.}
\]
The length constant is therefore the distance over which steady voltage falls to \(1/e\) of its starting value.

\paragraph{Concrete example.}
If \(\lambda=0.5\) mm, then a steady voltage deflection of \(10\) mV at a synapse produces, in the simplest passive semi-infinite approximation,
\[
v(0.5\ \mathrm{mm})=10e^{-1}\approx3.7\ \mathrm{mV},
\]
and
\[
v(1.0\ \mathrm{mm})=10e^{-2}\approx1.35\ \mathrm{mV}.
\]
If the dendrite radius is quadrupled while \(r_m\) and \(r_L\) are unchanged, then \(\lambda\) doubles because \(\lambda\propto\sqrt a\).

\paragraph{Computational neuroscience example.}
Compartmental neuron models approximate a dendritic tree by many small connected cylinders. Each compartment has a membrane equation, and neighboring compartments exchange axial current proportional to voltage differences. This discretizes the cable PDE into a large ODE system. The approximation becomes more accurate as compartments become electrically short relative to \(\lambda\), but more expensive computationally.

\paragraph{AI and machine-learning example.}
Cable theory has influenced neuromorphic and biologically inspired architectures through local filtering and spatial computation. Dendritic subunits can be viewed as analog processors that mix, attenuate, and nonlinearly transform inputs before they reach the soma. In machine learning, this motivates models with structured local computation, gated branches, and compartment-like hidden states. The mapping is useful when it specifies equations; it should not be reduced to the slogan ``dendrites are deep networks.''

\paragraph{Common misconceptions and failure modes.}
First, the cable equation is not the same as ordinary diffusion: membrane leakage contributes the \(-v\) term, so voltage both spreads and decays. Second, the length constant is not a fixed physical length of a dendrite; it depends on membrane and axial properties. Third, boundary conditions are not technical decoration. A sealed end and a voltage-clamped end can produce different voltage profiles. Fourth, a point-neuron model is justified only when the relevant cable is electrically compact for the signals being studied.

\paragraph{Exercises.}
\begin{enumerate}[label=\textbf{43.4.\arabic*.}, leftmargin=*]
\item \textbf{Mechanical.} If \(r_m\) is multiplied by \(4\) and all other cable parameters are fixed, by what factor do \(\tau_m\) and \(\lambda\) change?
\item \textbf{Proof or derivation.} Starting from \(\lambda^2v''-v=0\), derive the bounded semi-infinite solution \(v(x)=v(0)e^{-x/\lambda}\).
\item \textbf{Numerical/computational.} A passive voltage signal is \(8\) mV at \(x=0\), with \(\lambda=0.4\) mm. Compute the predicted signal at \(x=0.4\) mm and \(x=0.8\) mm.
\item \textbf{Interpretation.} Why can a distal synapse have a smaller somatic effect than a proximal synapse even if the local synaptic current is identical?
\item \textbf{Misconception/debugging.} A student simulates a finite cable but gets different answers after changing only the endpoint condition. Explain why this may be correct rather than a numerical bug.
\end{enumerate}

\paragraph{Closing bridge.}
This chapter built neurons from electrical conservation laws: capacitors, conductances, gating variables, and spatial cables. Chapter 44 now asks what can be simplified when the exact spike waveform is less important than spike timing, firing rate, phase-plane geometry, and population activity.

\clearpage
